\section{Conclusion}

This study examines whether a multimodal fact-checker benefits from knowing when evidence appeared and when several documents came from the same source. \systemname{} enforces a decision-time cutoff, groups dependent evidence before aggregation, checks whether the retrieved material is sufficient, and abstains when confidence is too low. The evaluation separates ordinary future-event performance from gains caused by temporal leakage and repeated documents. Its ablations also test whether the graph, source clusters, and sufficiency gate earn their place in the model. The final claims depend on controlled comparisons; a large stack of web documents does not count as corroboration on its own.

